% ============================================================================
%  ESTADÍSTICA DESCRIPTIVA — Sesión 8
%  Semana 8: 28 de septiembre al 3 de octubre de 2026
%  2.3 Medidas de dispersión: rango, varianza y desviación estándar
%  2.4 Medidas de forma: asimetría y curtosis
%  Duración prevista: 3 horas
% ============================================================================
\documentclass[aspectratio=169,10pt]{beamer}

\usepackage[utf8]{inputenc}
\IfFileExists{spanish.ldf}{\usepackage[spanish,es-noshorthands]{babel}}{}
\usepackage{amsmath,amssymb}
\usepackage{booktabs}
\usepackage{tikz}
\usepackage{listings}
\usepackage{pgfplots}
\pgfplotsset{compat=1.16}

\usetheme{Uniminuto}

\uniasignatura{Estadística descriptiva}
\uniprofesor{Andrés Felipe Puerta Velez}
\unicorreo{andres.puerta-v@uniminuto.edu.co}
\uninrc{CBASBL031 - CBASBL141}
\unisemana{Sesión 8 — Semana 8 · 28 de septiembre al 3 de octubre de 2026}

\newcommand{\listatiny}{\setbeamerfont{itemize/enumerate body}{size=\tiny}\tiny}

\pgfplotsset{
  uniplot/.style={
    width=8.0cm, height=3.5cm, axis lines=left,
    tick label style={font=\tiny,color=uniGris},
    label style={font=\tiny,color=uniGris},
    title style={font=\scriptsize\bfseries,color=uniAzul},
    every axis plot/.append style={fill=uniAzul!70,draw=uniAzul},
    ymajorgrids, grid style={uniGris!25,very thin},
  }
}

\newcommand{\fx}[1]{{\ttfamily\color{uniAzul}#1}}

% ------------------------------------------------------------- CÓDIGO R ----
\lstdefinestyle{Rcode}{
  language=R,
  basicstyle=\ttfamily\scriptsize,
  keywordstyle=\color{uniAzul}\bfseries,
  commentstyle=\color{uniGris}\itshape,
  stringstyle=\color{uniTitulo},
  showstringspaces=false, numbers=none, breaklines=true,
  columns=fullflexible,
  backgroundcolor=\color{uniCremaClaro},
  frame=leftline, framerule=1.2pt, rulecolor=\color{uniNaranja},
  xleftmargin=0.35cm, framexleftmargin=0.35cm,
  framexrightmargin=0.15cm, framextopmargin=2pt, framexbottommargin=2pt,
  aboveskip=0.15cm, belowskip=0.15cm
}
\lstdefinestyle{Rmini}{
  language=R,
  basicstyle=\ttfamily\tiny,
  keywordstyle=\color{uniAzul}\bfseries,
  commentstyle=\color{uniGris}\itshape,
  stringstyle=\color{uniTitulo},
  showstringspaces=false, numbers=none, breaklines=true,
  columns=fullflexible,
  backgroundcolor=\color{uniCremaClaro},
  frame=leftline, framerule=1.2pt, rulecolor=\color{uniNaranja},
  xleftmargin=0.35cm, framexleftmargin=0.35cm,
  framexrightmargin=0.15cm, framextopmargin=2pt, framexbottommargin=2pt,
  aboveskip=0.12cm, belowskip=0.12cm
}
\lstdefinestyle{Rout}{
  basicstyle=\ttfamily\tiny\color{black},
  showstringspaces=false, numbers=none, breaklines=true,
  columns=fullflexible,
  backgroundcolor=\color{uniGris!10},
  frame=leftline, framerule=1.2pt, rulecolor=\color{uniGris},
  xleftmargin=0.35cm, framexleftmargin=0.35cm,
  framexrightmargin=0.15cm, framextopmargin=2pt, framexbottommargin=2pt,
  aboveskip=0.12cm, belowskip=0.12cm
}

\begin{document}

\portada

\begin{frame}{Orden del día}
  \begin{enumerate}
    \item Llamado a lista y revisión de los ejercicios propuestos
    \item ¿Por qué el centro no basta?
    \item \azul{Tema 2.3: rango, varianza, desviación estándar y CV}
    \item \azul{Tema 2.4: asimetría y curtosis}
    \item Taller guiado con el caso Cafetería El Portal
    \item Errores frecuentes y cierre
  \end{enumerate}

  \vspace{0.25cm}
  {\footnotesize\color{uniGris} Sesión de 3 horas. Hoy se hacen muchas cuentas:
  traiga calculadora.}
\end{frame}

\begin{frame}{Por qué el centro no basta}
  \footnotesize
  \begin{columns}[T,onlytextwidth]
    \begin{column}{0.48\textwidth}
      \begin{block}{Sucursal A}
        \scriptsize
        Ventas diarias: \ttfamily 48 \ 49 \ 50 \ 51 \ 52 \rmfamily
        \[ \bar{x} = 50 \]
      \end{block}
    \end{column}
    \begin{column}{0.48\textwidth}
      \begin{block}{Sucursal B}
        \scriptsize
        Ventas diarias: \ttfamily 10 \ 30 \ 50 \ 70 \ 90 \rmfamily
        \[ \bar{x} = 50 \]
      \end{block}
    \end{column}
  \end{columns}

  \vspace{0.2cm}
  \begin{alertblock}{Misma media, negocios completamente distintos}
    \footnotesize
    La sucursal A es \azul{predecible}: se puede planear el inventario. La B es
    \azul{volátil}: unos días sobra mercancía y otros falta. Un solo número —la
    media— \destacar{no distingue estas dos situaciones}.
  \end{alertblock}

  \vspace{0.15cm}
  {\footnotesize\color{uniGris} Las medidas de dispersión responden a la
  pregunta: \azul{¿qué tan lejos del centro están los datos?}}
\end{frame}

\seccion{2.3 Medidas de dispersión}

\begin{frame}{Rango y rango intercuartílico}
  \footnotesize
  \begin{columns}[T,onlytextwidth]
    \begin{column}{0.48\textwidth}
      \begin{block}{Rango}
        \scriptsize
        \[ R = x_{\max} - x_{\min} \]
        Almacén Bello: $R = 352 - 112 = \destacar{240}$

        \vspace{0.1cm}
        \azul{Ventaja:} inmediato.
        \azul{Problema:} usa solo dos datos y los dos son los más extremos.
      \end{block}
    \end{column}
    \begin{column}{0.48\textwidth}
      \begin{exampleblock}{Rango intercuartílico}
        \scriptsize
        \[ RIC = Q_3 - Q_1 \]
        Almacén Bello: $267{,}56 - 186{,}29 = \destacar{81{,}27}$

        \vspace{0.1cm}
        Es el ancho del \azul{50\% central}. No lo afectan los valores
        extremos: es la medida robusta de dispersión.
      \end{exampleblock}
    \end{column}
  \end{columns}

  \vspace{0.15cm}
  \begin{alertblock}{Cómo se lee}
    \scriptsize
    ``El 50\% central de los clientes se mueve en una franja de
    \destacar{81 mil pesos} de ancho''. Es una frase mucho más útil que ``el
    rango es 240'', que solo describe a los dos clientes más raros.
  \end{alertblock}
\end{frame}

\begin{frame}{Varianza y desviación estándar — la idea}
  \footnotesize
  \begin{block}{El razonamiento, paso a paso}
    \scriptsize
    \begin{enumerate}\setlength{\itemsep}{-0.05cm}
      \item Queremos la distancia promedio de cada dato a la media:
            $x_i - \bar{x}$.
      \item Pero $\sum(x_i - \bar{x}) = 0$ siempre: los positivos cancelan a
            los negativos.
      \item Solución: \azul{elevar al cuadrado} antes de promediar. Eso es la
            \destacar{varianza}.
      \item La varianza queda en unidades al cuadrado (pesos$^2$), que no
            significan nada. Se saca la raíz: eso es la \destacar{desviación
            estándar}.
    \end{enumerate}
  \end{block}

  \vspace{0.1cm}
  \begin{columns}[T,onlytextwidth]
    \begin{column}{0.48\textwidth}
      \begin{exampleblock}{Muestra}
        \scriptsize
        \[ s^2 = \frac{\sum (x_i-\bar{x})^2}{n-1} \qquad s = \sqrt{s^2} \]
      \end{exampleblock}
    \end{column}
    \begin{column}{0.48\textwidth}
      \begin{exampleblock}{Población}
        \scriptsize
        \[ \sigma^2 = \frac{\sum (x_i-\mu)^2}{N} \qquad \sigma = \sqrt{\sigma^2} \]
      \end{exampleblock}
    \end{column}
  \end{columns}

  \vspace{0.05cm}
  {\tiny\color{uniGris} La diferencia entre $n-1$ y $N$ importa poco cuando $n$
  es grande, pero hay que declarar cuál se usó. En este curso, si el enunciado
  dice ``muestra'', se divide entre $n-1$.}
\end{frame}

\begin{frame}{Varianza con datos agrupados — el procedimiento}
  \scriptsize
  \[ s^2 = \frac{\sum f_i (x_i-\bar{x})^2}{n-1} \qquad\qquad
     \sigma^2 = \frac{\sum f_i (x_i-\bar{x})^2}{n} \]

  \vspace{0.1cm}
  \begin{center}
  \begin{tabular}{@{}lrrrrr@{}}
    \toprule
    \textbf{Clase} & \textbf{$x_i$} & \textbf{$f_i$} &
    \textbf{$x_i-\bar{x}$} & \textbf{$(x_i-\bar{x})^2$} &
    \textbf{$f_i(x_i-\bar{x})^2$} \\
    \midrule
    $[112,152)$ & 132 &  4 & $-95$ &  9\,025 &  36\,100 \\
    $[152,192)$ & 172 &  7 & $-55$ &  3\,025 &  21\,175 \\
    $[192,232)$ & 212 & 11 & $-15$ &     225 &   2\,475 \\
    $[232,272)$ & 252 &  9 & $+25$ &     625 &   5\,625 \\
    $[272,312)$ & 292 &  6 & $+65$ &  4\,225 &  25\,350 \\
    $[312,352]$ & 332 &  3 & $+105$ & 11\,025 & 33\,075 \\
    \midrule
    \textbf{Total} & & \textbf{40} & & & \textbf{123\,800} \\
    \bottomrule
  \end{tabular}
  \end{center}

  \vspace{0.05cm}
  {\tiny\color{uniGris} Recuerde de la sesión 7: $\bar{x} = 227{,}0$. Las
  desviaciones se calculan respecto a esa media, no respecto a cada clase.}
\end{frame}

\begin{frame}{Varianza y desviación del Almacén Bello}
  \footnotesize
  \begin{columns}[T,onlytextwidth]
    \begin{column}{0.48\textwidth}
      \begin{exampleblock}{El cálculo}
        \scriptsize
        \[ s^2 = \frac{123\,800}{40-1} = \destacar{3\,174{,}36} \]
        \[ s = \sqrt{3\,174{,}36} = \destacar{56{,}34} \]
        En unidades: pesos$^2$ la varianza, \azul{miles de pesos} la
        desviación.
      \end{exampleblock}
    \end{column}
    \begin{column}{0.48\textwidth}
      \begin{block}{Cómo se interpreta}
        \scriptsize
        El gasto de un cliente se aparta del promedio, \azul{en promedio}, unos
        \destacar{56\,340 pesos}.

        \vspace{0.1cm}
        Con datos originales: $s = 57{,}71$. La diferencia es, otra vez, el
        error de agrupamiento.
      \end{block}
    \end{column}
  \end{columns}

  \vspace{0.1cm}
  \begin{alertblock}{La regla empírica (distribuciones aproximadamente simétricas)}
    \scriptsize
    Cerca del \destacar{68\%} de los datos cae en $\bar{x} \pm s$;
    el \destacar{95\%} en $\bar{x} \pm 2s$; casi todo en $\bar{x} \pm 3s$.
    Aquí: $227 \pm 56{,}34 = [170{,}7,\ 283{,}3]$, franja que en efecto
    contiene a la mayoría de los clientes.
  \end{alertblock}
\end{frame}

\begin{frame}{Coeficiente de variación — comparar lo incomparable}
  \footnotesize
  \begin{columns}[T,onlytextwidth]
    \begin{column}{0.38\textwidth}
      \begin{block}{Fórmula}
        \scriptsize
        \[ CV = \frac{s}{\bar{x}} \times 100 \]
        Se expresa en porcentaje y \azul{no tiene unidades}.
        \vspace{0.1cm}
        Almacén Bello:
        \[ CV = \frac{56{,}34}{227{,}0}\times 100 = \destacar{24{,}8\%} \]
      \end{block}
    \end{column}
    \begin{column}{0.58\textwidth}
      \begin{exampleblock}{Para qué sirve: dos productos}
        \scriptsize
        \begin{tabular}{@{}lrrr@{}}
          \toprule
          & \textbf{$\bar{x}$} & \textbf{$s$} & \textbf{$CV$} \\
          \midrule
          Producto A (unidades/día) & 120 & 18 & 15{,}0\% \\
          Producto B (unidades/día) &  45 &  9 & 20{,}0\% \\
          \bottomrule
        \end{tabular}
        \vspace{0.1cm}
        A tiene una desviación \azul{mayor} (18 frente a 9), pero
        \destacar{B es el producto más inestable} en términos relativos.
      \end{exampleblock}
    \end{column}
  \end{columns}

  \vspace{0.1cm}
  \begin{alertblock}{Guía de lectura (orientativa)}
    \scriptsize
    $CV < 10\%$: muy homogéneo. \quad $10$–$25\%$: homogeneidad
    aceptable. \quad $25$–$50\%$: heterogéneo. \quad $> 50\%$: muy
    disperso, la media deja de ser un buen resumen.
    \azul{Solo tiene sentido si $\bar{x} > 0$.}
  \end{alertblock}
\end{frame}

\begin{frame}{Ejercicio 1 — en clase}
  \begin{alertblock}{En parejas (15 minutos)}
    \footnotesize
    Dos cajeros registran el tiempo de atención (minutos) de 5 clientes:

    \vspace{0.15cm}
    \centering\scriptsize
    \begin{tabular}{@{}ll@{}}
      \toprule
      \textbf{Cajero 1} & 4 \quad 5 \quad 6 \quad 7 \quad 8 \\
      \textbf{Cajero 2} & 2 \quad 3 \quad 6 \quad 9 \quad 10 \\
      \bottomrule
    \end{tabular}

    \vspace{0.2cm}
    \raggedright
    \begin{enumerate}\setlength{\itemsep}{-0.05cm}
      \item Calcule la media de cada cajero.
      \item Calcule el rango de cada uno.
      \item Calcule $s^2$ y $s$ de cada uno (como muestra, $n-1$).
      \item Calcule el $CV$ de cada uno.
      \item ¿Cuál cajero es más consistente? ¿Con qué argumento lo defendería
            ante el jefe?
    \end{enumerate}
  \end{alertblock}
\end{frame}

\begin{frame}{Ejercicio 1 — solución}
  \scriptsize
  \begin{columns}[T,onlytextwidth]
    \begin{column}{0.49\textwidth}
      \begin{block}{Cajero 1}
        \scriptsize
        $\bar{x} = 30/5 = 6{,}0$ \quad $R = 4$
        \[ \sum(x-\bar{x})^2 = 4+1+0+1+4 = 10 \]
        \[ s^2 = \tfrac{10}{4} = 2{,}5 \qquad s = \destacar{1{,}58} \]
        \[ CV = \tfrac{1{,}58}{6{,}0}\times100 = \destacar{26{,}4\%} \]
      \end{block}
    \end{column}
    \begin{column}{0.49\textwidth}
      \begin{block}{Cajero 2}
        \scriptsize
        $\bar{x} = 30/5 = 6{,}0$ \quad $R = 8$
        \[ \sum(x-\bar{x})^2 = 16+9+0+9+16 = 50 \]
        \[ s^2 = \tfrac{50}{4} = 12{,}5 \qquad s = \destacar{3{,}54} \]
        \[ CV = \tfrac{3{,}54}{6{,}0}\times100 = \destacar{58{,}9\%} \]
      \end{block}
    \end{column}
  \end{columns}

  \vspace{0.1cm}
  \begin{alertblock}{Punto 5}
    \scriptsize
    Los dos atienden en \azul{6 minutos promedio}, pero el cajero 1 tiene una
    desviación de 1,6 minutos frente a 3,5. \destacar{El cajero 1 es más
    consistente}: su cliente puede prever cuánto va a esperar. Al jefe se le
    dice con el CV: 26\% frente a 59\%.
  \end{alertblock}
\end{frame}

\seccion{2.4 Medidas de forma}

\begin{frame}{Asimetría — qué mide}
  \centering
  \begin{tikzpicture}[scale=0.82]
    \begin{scope}
      \foreach \x/\h in {0/0.4,0.5/0.9,1/1.5,1.5/1.9,2/1.5,2.5/0.9,3/0.4}
        \fill[uniAzul!70] (\x,0) rectangle (\x+0.5,\h);
      \draw[uniGris] (0,0)--(3.5,0);
      \node[font=\tiny,color=uniAzul,align=center] at (1.75,-0.45)
        {Simétrica\\$As = 0$};
    \end{scope}
    \begin{scope}[xshift=5.0cm]
      \foreach \x/\h in {0/1.9,0.5/1.6,1/1.1,1.5/0.8,2/0.5,2.5/0.3,3/0.15}
        \fill[uniNaranja!85] (\x,0) rectangle (\x+0.5,\h);
      \draw[uniGris] (0,0)--(3.5,0);
      \node[font=\tiny,color=uniAzul,align=center] at (1.75,-0.45)
        {Positiva (a la derecha)\\$As > 0$ \ · \ $Mo<Me<\bar{x}$};
    \end{scope}
    \begin{scope}[xshift=10.0cm]
      \foreach \x/\h in {0/0.15,0.5/0.3,1/0.5,1.5/0.8,2/1.1,2.5/1.6,3/1.9}
        \fill[uniNaranja!85] (\x,0) rectangle (\x+0.5,\h);
      \draw[uniGris] (0,0)--(3.5,0);
      \node[font=\tiny,color=uniAzul,align=center] at (1.75,-0.45)
        {Negativa (a la izquierda)\\$As < 0$ \ · \ $\bar{x}<Me<Mo$};
    \end{scope}
  \end{tikzpicture}

  \vspace{0.5cm}
  \begin{alertblock}{La asimetría es una pregunta de negocio}
    \scriptsize
    Gastos, salarios y tiempos de espera casi siempre son
    \destacar{asimétricos positivos}: hay un piso pero no un techo. Las notas
    de un curso fácil suelen ser \destacar{asimétricas negativas}: casi todos
    arriba y unos pocos muy abajo.
  \end{alertblock}
\end{frame}

\begin{frame}{Los tres coeficientes de asimetría}
  \scriptsize
  \begin{center}
  \begin{tabular}{@{}p{3.3cm}p{4.0cm}p{4.6cm}@{}}
    \toprule
    \textbf{Coeficiente} & \textbf{Fórmula} & \textbf{Cuándo usarlo} \\
    \midrule
    Pearson (moda) &
      $As_1 = \dfrac{\bar{x}-Mo}{s}$ &
      Rápido; exige que la moda esté bien definida \\
    \addlinespace
    Pearson (mediana) &
      $As_2 = \dfrac{3(\bar{x}-Me)}{s}$ &
      El más usado en la práctica \\
    \addlinespace
    Bowley (cuartiles) &
      $As_B = \dfrac{Q_3+Q_1-2Q_2}{Q_3-Q_1}$ &
      Robusto: no lo afectan los extremos \\
    \addlinespace
    Fisher (momentos) &
      $g_1 = \dfrac{\sum f_i(x_i-\bar{x})^3 / n}{\sigma^3}$ &
      El que reportan R y los paquetes estadísticos \\
    \bottomrule
  \end{tabular}
  \end{center}

  \vspace{0.1cm}
  \begin{alertblock}{Lo único que hay que leer del resultado}
    \scriptsize
    El \destacar{signo} y la \destacar{magnitud}: positivo, cola a la derecha;
    negativo, cola a la izquierda; cerca de cero, simétrica. Valores entre
    $-0{,}5$ y $0{,}5$ se consideran aproximadamente simétricos.
  \end{alertblock}
\end{frame}

\begin{frame}{Asimetría del Almacén Bello}
  \scriptsize
  \begin{columns}[T,onlytextwidth]
    \begin{column}{0.49\textwidth}
      \begin{exampleblock}{Los datos que ya tenemos}
        \scriptsize
        \[ \bar{x}=227{,}0 \quad Me=224{,}73 \quad Mo=218{,}67 \]
        \[ s = 56{,}34 \quad Q_1=186{,}29 \quad Q_3=267{,}56 \]
      \end{exampleblock}
    \end{column}
    \begin{column}{0.49\textwidth}
      \begin{block}{Los tres cálculos}
        \scriptsize
        \[ As_1 = \frac{227{,}0-218{,}67}{56{,}34} = \destacar{0{,}148} \]
        \[ As_2 = \frac{3(227{,}0-224{,}73)}{56{,}34} = \destacar{0{,}121} \]
        \[ As_B = \frac{267{,}56+186{,}29-2(224{,}73)}{81{,}27}
                = \destacar{0{,}054} \]
      \end{block}
    \end{column}
  \end{columns}

  \vspace{0.1cm}
  \begin{alertblock}{La conclusión}
    \scriptsize
    Los tres coeficientes son \destacar{positivos pero pequeños}: hay una cola
    a la derecha —los clientes de gasto alto— pero la distribución no está muy
    deformada. Coincide con lo que muestra el histograma y con el orden
    $Mo < Me < \bar{x}$.
  \end{alertblock}
\end{frame}

\begin{frame}{Curtosis — qué tan concentrada está la distribución}
  \centering
  \begin{tikzpicture}[scale=0.9]
    \draw[uniGris,->] (-3.2,0)--(3.2,0);
    \draw[uniAzul,very thick,domain=-3:3,samples=80]
      plot (\x,{2.0*exp(-\x*\x*1.6)});
    \draw[uniTitulo,very thick,dashed,domain=-3:3,samples=80]
      plot (\x,{1.15*exp(-\x*\x*0.42)});
    \node[font=\tiny,color=uniAzul,right] at (0.9,1.75) {Leptocúrtica ($g_2>0$)};
    \node[font=\tiny,color=uniTitulo,right] at (1.5,0.85) {Platicúrtica ($g_2<0$)};
  \end{tikzpicture}

  \vspace{0.2cm}
  \begin{columns}[T,onlytextwidth]
    \begin{column}{0.46\textwidth}
      \begin{block}{Fórmula (exceso de curtosis)}
        \scriptsize
        \[ g_2 = \frac{\sum f_i (x_i-\bar{x})^4 / n}{\sigma^4} - 3 \]
      \end{block}
    \end{column}
    \begin{column}{0.50\textwidth}
      \begin{exampleblock}{Cómo se lee}
        \scriptsize
        $g_2 > 0$: \azul{leptocúrtica}, pico alto y colas pesadas. \\
        $g_2 = 0$: \azul{mesocúrtica}, como la normal. \\
        $g_2 < 0$: \azul{platicúrtica}, achatada, datos repartidos.
      \end{exampleblock}
    \end{column}
  \end{columns}
\end{frame}

\begin{frame}{Curtosis del Almacén Bello}
  \scriptsize
  \begin{columns}[T,onlytextwidth]
    \begin{column}{0.52\textwidth}
      \centering
      \begin{tabular}{@{}lrrr@{}}
        \toprule
        \textbf{$x_i$} & \textbf{$f_i$} & \textbf{$f_i d_i^3$} &
        \textbf{$f_i d_i^4$} \\
        \midrule
        132 &  4 & $-3\,429\,500$ & 325\,802\,500 \\
        172 &  7 & $-1\,164\,625$ &  64\,054\,375 \\
        212 & 11 & $\ \ \ -37\,125$ &     556\,875 \\
        252 &  9 & $+140\,625$ &   3\,515\,625 \\
        292 &  6 & $+1\,647\,750$ & 107\,103\,750 \\
        332 &  3 & $+3\,472\,875$ & 364\,651\,875 \\
        \midrule
        \textbf{Suma} & \textbf{40} & \textbf{630\,000} &
        \textbf{865\,685\,000} \\
        \bottomrule
      \end{tabular}
      {\tiny $d_i = x_i - 227{,}0$}
    \end{column}
    \begin{column}{0.46\textwidth}
      \begin{block}{Los coeficientes de Fisher}
        \scriptsize
        Con $\sigma = 55{,}63$ (dividiendo entre $n$):
        \[ g_1 = \frac{630\,000/40}{55{,}63^3} = \destacar{0{,}091} \]
        \[ g_2 = \frac{865\,685\,000/40}{55{,}63^4} - 3 \]
        \[ g_2 = 2{,}259 - 3 = \destacar{-0{,}741} \]
      \end{block}
    \end{column}
  \end{columns}

  \vspace{0.05cm}
  \begin{alertblock}{Lectura conjunta}
    \tiny
    $g_1 = 0{,}09$: prácticamente simétrica según Fisher. \quad
    $g_2 = -0{,}74$: \destacar{platicúrtica}, más achatada que la normal. Los
    gastos están repartidos en una franja ancha, sin una concentración fuerte
    en un punto.
  \end{alertblock}
\end{frame}

\begin{frame}{Ejercicio 2 — en clase (taller principal)}
  \begin{alertblock}{En grupos de 3 (30 minutos)}
    \scriptsize
    Cafetería El Portal ($n = 45$). De la sesión 7 ya sabemos:
    $\bar{x} = 22{,}33$, \ $Me = 21{,}92$, \ $Mo = 21{,}43$, \
    $Q_1 = 16{,}17$, \ $Q_3 = 28{,}05$.

    \vspace{0.15cm}
    \centering
    \begin{tabular}{@{}lrrrrrr@{}}
      \toprule
      \textbf{Clase} & $[6,12)$ & $[12,18)$ & $[18,24)$ & $[24,30)$ &
      $[30,36)$ & $[36,42]$ \\
      \midrule
      $x_i$ &  9 & 15 & 21 & 27 & 33 & 39 \\
      $f_i$ &  5 &  9 & 13 & 10 &  5 &  3 \\
      \bottomrule
    \end{tabular}

    \vspace{0.2cm}
    \raggedright
    \begin{enumerate}\setlength{\itemsep}{-0.07cm}
      \item Construya la columna $f_i(x_i-\bar{x})^2$ y súmela.
      \item Calcule $s^2$ y $s$ (como muestra).
      \item Calcule el $CV$ e interprete la homogeneidad del grupo.
      \item Calcule el $RIC$.
      \item Calcule $As_1$, $As_2$ y $As_B$. ¿Coinciden en el signo?
      \item Compare el $CV$ de El Portal con el del Almacén Bello (24,8\%).
            ¿Qué negocio tiene clientes más parecidos entre sí?
    \end{enumerate}
  \end{alertblock}
\end{frame}

\begin{frame}{Ejercicio 2 — solución (dispersión)}
  \scriptsize
  \begin{columns}[T,onlytextwidth]
    \begin{column}{0.50\textwidth}
      \centering
      \begin{tabular}{@{}rrrr@{}}
        \toprule
        \textbf{$x_i$} & \textbf{$f_i$} & \textbf{$(x_i-\bar{x})^2$} &
        \textbf{$f_i(x_i-\bar{x})^2$} \\
        \midrule
         9 &  5 & 177{,}78 & 888{,}89 \\
        15 &  9 &  53{,}78 & 484{,}00 \\
        21 & 13 &   1{,}78 &  23{,}11 \\
        27 & 10 &  21{,}78 & 217{,}78 \\
        33 &  5 & 113{,}78 & 568{,}89 \\
        39 &  3 & 277{,}78 & 833{,}33 \\
        \midrule
        & \textbf{45} & & \textbf{3\,016{,}00} \\
        \bottomrule
      \end{tabular}
    \end{column}
    \begin{column}{0.47\textwidth}
      \begin{exampleblock}{Puntos 2 a 4}
        \scriptsize
        \[ s^2 = \frac{3\,016}{44} = \destacar{68{,}55} \]
        \[ s = \destacar{8{,}28} \text{ miles de pesos} \]
        \[ CV = \frac{8{,}28}{22{,}33}\times100 = \destacar{37{,}1\%} \]
        \[ RIC = 28{,}05-16{,}17 = \destacar{11{,}88} \]
      \end{exampleblock}
    \end{column}
  \end{columns}

  \vspace{0.05cm}
  {\tiny\color{uniGris} Con $CV = 37{,}1\%$ el grupo es \azul{heterogéneo}:
  el consumo varía bastante de un cliente a otro.}
\end{frame}

\begin{frame}{Ejercicio 2 — solución (forma y comparación)}
  \scriptsize
  \begin{columns}[T,onlytextwidth]
    \begin{column}{0.50\textwidth}
      \begin{block}{Punto 5: los tres coeficientes}
        \scriptsize
        \[ As_1 = \frac{22{,}33-21{,}43}{8{,}28} = \destacar{0{,}109} \]
        \[ As_2 = \frac{3(22{,}33-21{,}92)}{8{,}28} = \destacar{0{,}149} \]
        \[ As_B = \frac{28{,}05+16{,}17-2(21{,}92)}{11{,}88}
                = \destacar{0{,}031} \]
        Los tres son \azul{positivos}: coinciden.
      \end{block}
    \end{column}
    \begin{column}{0.47\textwidth}
      \begin{exampleblock}{Punto 6: la comparación}
        \scriptsize
        \begin{tabular}{@{}lrr@{}}
          \toprule
          & \textbf{Bello} & \textbf{El Portal} \\
          \midrule
          $\bar{x}$ & 227{,}0 & 22{,}33 \\
          $s$       &  56{,}34 &  8{,}28 \\
          $CV$      & 24{,}8\% & 37{,}1\% \\
          \bottomrule
        \end{tabular}
        \vspace{0.1cm}
        El \destacar{Almacén Bello} tiene la clientela más homogénea, aunque
        su desviación en pesos sea mucho mayor.
      \end{exampleblock}
    \end{column}
  \end{columns}

  \vspace{0.05cm}
  {\tiny\color{uniGris} Esta es la razón de ser del CV: comparar dispersión
  entre variables con medias y unidades muy distintas.}
\end{frame}

\begin{frame}{Errores frecuentes}
  \scriptsize
  \begin{center}
  \begin{tabular}{@{}p{5.4cm}p{6.4cm}@{}}
    \toprule
    \textbf{Error} & \textbf{Corrección} \\
    \midrule
    Reportar la varianza como si describiera los datos &
      La varianza está en unidades al cuadrado; se reporta $s$ \\
    Dividir entre $n$ cuando el enunciado dice ``muestra'' &
      Muestra: $n-1$. Población: $N$. Declararlo siempre \\
    Comparar desviaciones de variables con medias distintas &
      Para eso está el $CV$ \\
    Calcular $CV$ con media negativa o cero &
      No tiene interpretación; no se usa \\
    Confundir el signo de la asimetría con ``bueno o malo'' &
      Solo indica hacia dónde está la cola \\
    Olvidar restar 3 en la curtosis de Fisher &
      $g_2 = m_4/\sigma^4 - 3$; sin el $-3$ se compara contra 3 \\
    Redondear la media antes de calcular las desviaciones &
      Arrastra el error a toda la columna \\
    \bottomrule
  \end{tabular}
  \end{center}
\end{frame}

\begin{frame}[fragile]{Dispersión y forma, en R}
  \tiny
\begin{lstlisting}[style=Rmini]
diff(range(gasto))              # rango
IQR(gasto)                      # rango intercuartilico
var(gasto);  sd(gasto)          # varianza y desviacion DE MUESTRA (n-1)
sd(gasto) / mean(gasto) * 100   # coeficiente de variacion

# R base no trae asimetria ni curtosis: se calculan a mano
n  <- length(gasto);  mu <- mean(gasto)
sg <- sqrt(sum((gasto - mu)^2) / n)        # sigma poblacional
As2 <- 3 * (mu - median(gasto)) / sd(gasto)
g1  <- sum((gasto - mu)^3) / n / sg^3
g2  <- sum((gasto - mu)^4) / n / sg^4 - 3
round(c(As2 = As2, g1 = g1, g2 = g2), 3)
\end{lstlisting}
\begin{lstlisting}[style=Rout]
[1] 240
[1] 76.75
[1] 3329.977
[1] 57.706
[1] 25.404
   As2     g1     g2
-0.018  0.077 -0.564
\end{lstlisting}

  \vspace{-0.28cm}
  \begin{alertblock}{Una lección escondida en estos números}
    \tiny
    Con datos crudos $As_2=-0{,}018$ y con la tabla agrupada dio $+0{,}121$:
    cuando la asimetría es \destacar{muy leve}, los coeficientes basados en la
    mediana cambian de signo con nada. El de Fisher ($g_1=+0{,}077$) y el
    histograma coinciden: cola derecha, pero pequeña.
  \end{alertblock}
\end{frame}

\begin{frame}{Cierre de la sesión}
  \begin{block}{Lo que debe quedar claro hoy}
    \scriptsize
    \begin{itemize}\setlength{\itemsep}{0.02cm}
      \item Dos grupos con la misma media pueden ser negocios distintos: la
            \azul{dispersión} lo revela.
      \item $R$ y $RIC$ son rápidos; $s^2$ y $s$ usan toda la información.
      \item $s$ se reporta en las \azul{unidades originales}; la varianza no.
      \item El $CV$ permite \azul{comparar} dispersión entre variables
            distintas.
      \item $As$ indica hacia dónde está la cola; $g_2$, qué tan apuntada es
            la distribución.
    \end{itemize}
  \end{block}

  \vspace{0.1cm}
  \begin{alertblock}{Trabajo independiente}
    \scriptsize
    Resolver los 20 ejercicios de la lámina siguiente. En la sesión 9
    construimos el \destacar{diagrama de caja}, que reúne en un solo dibujo la
    mediana, los cuartiles, el recorrido y los valores atípicos.
  \end{alertblock}
\end{frame}

\begin{frame}{Ejercicios propuestos}
  \vspace{-0.1cm}
  \begin{beamercolorbox}[wd=\textwidth,sep=0.10cm]{block body}
    \centering\scriptsize
    \textbf{1--8:} cálculo directo. \quad \textbf{9--14:} agrupados. \quad
    \textbf{15--20:} interpretación y comparación.
  \end{beamercolorbox}
  \listatiny
  \begin{columns}[T,onlytextwidth]
    \begin{column}{0.47\textwidth}
      \begin{enumerate}\setlength{\itemsep}{0.16cm}
        \item Rango de: 12, 19, 7, 25, 14
        \item Media de: 2, 4, 4, 4, 5, 5, 7, 9
        \item $\sum(x-\bar{x})^2$ de ese conjunto
        \item $s^2$ y $s$ como muestra
        \item $\sigma^2$ y $\sigma$ como población
        \item $CV$ de ese conjunto
        \item Si $s = 12$ y $\bar{x} = 60$: ¿$CV$?
        \item Si $CV = 30\%$ y $\bar{x} = 80$: ¿$s$?
        \item $RIC$ del Almacén Bello
        \item $s$ agrupada del Almacén Bello (verifique 56{,}34)
      \end{enumerate}
    \end{column}
    \begin{column}{0.50\textwidth}
      \begin{enumerate}\setlength{\itemsep}{0.16cm}
        \setcounter{enumi}{10}
        \item $CV$ del Almacén Bello
        \item $As_2$ del Almacén Bello
        \item $As_B$ del Almacén Bello
        \item $g_2$ del Almacén Bello y su nombre
        \item Dos sucursales: $\bar{x}=200$, $s=20$ y $\bar{x}=50$, $s=8$.
              ¿Cuál es más estable?
        \item ¿Por qué se eleva al cuadrado en la varianza?
        \item ¿Qué significa $s = 0$?
        \item Si $As_2 < 0$, ¿dónde está la cola?
        \item ¿Qué medida de dispersión no se afecta por atípicos?
        \item ¿Por qué la varianza no se reporta en el informe final?
      \end{enumerate}
    \end{column}
  \end{columns}
\end{frame}

\begin{frame}{Respuestas}{Autoevaluación de los ejercicios propuestos}
  \listatiny
  \begin{columns}[T,onlytextwidth]
    \begin{column}{0.47\textwidth}
      \begin{enumerate}\setlength{\itemsep}{0.16cm}
        \item $25-7 = 18$
        \item $40/8 = 5{,}0$
        \item $9+1+1+1+0+0+4+16 = 32$
        \item $s^2 = 32/7 = 4{,}57$; \ $s = 2{,}14$
        \item $\sigma^2 = 32/8 = 4{,}00$; \ $\sigma = 2{,}00$
        \item $CV = 2{,}14/5{,}0 = 42{,}8\%$
        \item $20{,}0\%$
        \item $s = 0{,}30 \times 80 = 24$
        \item $267{,}56-186{,}29 = 81{,}27$
        \item $\sqrt{123\,800/39} = 56{,}34$
      \end{enumerate}
    \end{column}
    \begin{column}{0.50\textwidth}
      \begin{enumerate}\setlength{\itemsep}{0.16cm}
        \setcounter{enumi}{10}
        \item $56{,}34/227{,}0 = 24{,}8\%$
        \item $3(227{,}0-224{,}73)/56{,}34 = 0{,}121$
        \item $0{,}054$
        \item $-0{,}741$: platicúrtica
        \item $CV$: 10\% frente a 16\%; la primera es más estable
        \item Para que las desviaciones no se cancelen
        \item Todos los datos son iguales: no hay dispersión
        \item A la izquierda (asimetría negativa)
        \item El $RIC$ (y el $As_B$, en forma)
        \item Porque sus unidades están al cuadrado y no se interpretan
      \end{enumerate}
    \end{column}
  \end{columns}
\end{frame}

\begin{frame}{Referencias}
  \footnotesize
  \begin{itemize}
    \item[{[1]}] Monroy Saldívar, S. (2008). \textit{Estadística descriptiva}.
      Instituto Politécnico Nacional.
      \texttt{https://elibro.net/es/lc/uniminuto/titulos/74722}
    \item[{[2]}] Ordóñez Fernández, F. F. \& González Fernández, J. (2021).
      \textit{Estadística descriptiva paso a paso}. Ediciones Pirámide.
      \texttt{https://elibro.net/es/lc/uniminuto/titulos/215449}
    \item[{[3]}] Llinás Solano, H. (2014). \textit{Introducción a la
      estadística con aplicaciones en ciencias sociales}. Universidad del
      Norte. \texttt{https://elibro.net/es/lc/uniminuto/titulos/70062}
  \end{itemize}

  \vspace{0.35cm}
  \begin{center}
    {\Large\color{uniAzul}\textbf{¿Preguntas?}}
  \end{center}
\end{frame}

\end{document}
